\chapter{Traceability und Änderungsmanagement}

\begin{lernziele}
Nach diesem Kapitel können Sie eine durchgängige Nachvollziehbarkeit von Anforderungen bis Tests aufbauen.
\end{lernziele}

\section{Was ist Traceability?}
\gls{traceability} bedeutet Nachvollziehbarkeit. Sie zeigt, wie Artefakte zusammenhängen. Bei MBSE ist das einer der größten praktischen Vorteile.

\section{Kette im Roboterprojekt}
\begin{verbatim}
REQ-002 Hindernisse erkennen
-> PerceptionSystem
-> ObstacleDetector
-> obstacle_detector_node
-> /scan
-> TC-002 Hinderniserkennung
\end{verbatim}

\section{Nutzen}
Wenn sich REQ-002 ändert, kann analysiert werden, welche Komponenten betroffen sind. Vielleicht muss der Lidar ausgetauscht werden. Vielleicht muss der Test angepasst werden. Vielleicht reicht eine Softwareänderung.

\section{Impact Analysis}
Impact Analysis bedeutet Auswirkungsanalyse. Sie beantwortet die Frage: Was ändert sich, wenn eine Anforderung, Komponente oder Schnittstelle geändert wird?

\section{Typische Fehler}
Traceability wird oft zu spät aufgebaut. Dann ist sie mühsam. Besser ist es, Beziehungen während der Modellierung schrittweise zu pflegen.

\begin{uebungbox}
Wählen Sie eine Anforderung aus und erstellen Sie eine vollständige Traceability-Kette bis zu einem Testfall.
\end{uebungbox}
